\documentclass[a4paper,11pt,oneside]{article}
\usepackage[T1]{fontenc}
\usepackage[utf8]{inputenc}
\usepackage[english]{babel}
\usepackage{lmodern,microtype}
\usepackage{geometry,fancyhdr,graphicx,amsmath,amssymb,booktabs,tabularx,longtable,array}
\usepackage{enumitem,listings,xcolor,caption,placeins}
\usepackage[hidelinks]{hyperref}
\usepackage{xurl}
\geometry{left=24mm,right=24mm,top=28mm,bottom=25mm,headheight=26pt,headsep=10pt}
\graphicspath{{figures/}{./}}
\pagestyle{fancy}
\fancyhf{}
\fancyhead[L]{\small SeismicTorch: waveform classification}
\fancyhead[R]{\small Gauss Chang\\2026 project revision}
\fancyfoot[C]{\thepage}
\setlength{\parindent}{0pt}
\setlength{\parskip}{5pt}
\linespread{1.10}
\setlist{nosep,leftmargin=1.6em}
\captionsetup{font=small,labelfont=bf,skip=7pt}
\definecolor{accent}{HTML}{2475A7}
\definecolor{codebg}{HTML}{F1F4F7}
\lstset{basicstyle=\footnotesize\ttfamily,backgroundcolor=\color{codebg},frame=single,
  rulecolor=\color{black!15},breaklines=true,columns=fullflexible,keepspaces=true,
  showstringspaces=false,aboveskip=9pt,belowskip=9pt}
\newcommand{\code}[1]{\nolinkurl{#1}}
\newcommand{\fig}[5][]{\begin{figure}[htbp]\centering\includegraphics[width=#3\linewidth]{#2.pdf}\caption[#1]{#4}\label{#5}\end{figure}}
\newcolumntype{Y}{>{\raggedright\arraybackslash}X}
\makeatletter
\renewcommand{\l@section}[2]{\@dottedtocline{1}{0em}{1.8em}{\bfseries #1}{\bfseries #2}}
\makeatother
\hypersetup{pdftitle={Deep Learning Based Seismic Wave Identification},pdfauthor={Gauss Chang},pdfsubject={Verified QSIS waveform CNN, training diagnostics and standalone C++ inference}}

\begin{document}
\hypersetup{pageanchor=false}
\begin{titlepage}
\centering
\vspace*{5mm}
\includegraphics[width=.27\textwidth]{IES_logo.pdf}\\[9mm]
{\Large Institute of Earth Sciences\\Academia Sinica\par}
\vspace{8mm}
{\large Summer Internship Project / 2026 Technical Revision\par}
\vspace{12mm}
\rule{\linewidth}{1.2pt}\\[4mm]
{\huge\bfseries Deep Learning Based\\[2mm] Seismic Wave Identification\par}
\vspace{4mm}
\rule{\linewidth}{1.2pt}\\[7mm]
{\large Reproducible waveform classification, independent QSIS confirmation,\\and standalone embedded-model inference\par}
\vspace{17mm}
\begin{minipage}[t]{.43\textwidth}\raggedright
\textit{Author}\\Gauss Chang\\Department of Physics, NTU
\end{minipage}\hfill
\begin{minipage}[t]{.49\textwidth}\raggedleft
\textit{Advisors, original project}\\Dr. Kuo-Fong Ma\\Dr. Wen-Tzong Liang
\end{minipage}
\vfill
{\large 21 November 2026\par}
\vspace{5mm}
{\small Software release: \href{https://github.com/Gauss963/SeismicTorch/releases/tag/v1.0.0}{SeismicTorch v1.0.0}\\
Evidence-based revision of the original project report.}
\end{titlepage}
\hypersetup{pageanchor=true}

\pagenumbering{roman}
\begin{abstract}
This project classifies short three-component seismic waveforms into P waves,
S waves and noise, with an emphasis on QSIS MEMS acceleration recordings in
Taiwanese buildings. The revived implementation retains a compact waveform
CNN with convolutional channels $16\rightarrow32\rightarrow64$ and a classifier
$384\rightarrow64\rightarrow16\rightarrow3$. It contains 33,795 trainable
parameters and consumes two-second, 40 Hz waveform windows; spectrograms are
used only for diagnostics. The work includes GPU training on Apple MPS,
shared preprocessing, recording/event-aware evaluation, model-blind waveform
reference review and dependency-light native deployment.

The initial prepared-archive experiment reached 94.14\% on the STEAD test
archive but only 70.97\% on the QSIS test partition. These results motivated
arrival-window review, Taiwan strong-motion adaptation, training-only hard
background examples, data augmentation and per-record gravity alignment.
The released model was fitted for 20 epochs on 226 source windows initialized
from a six-epoch CWA adaptation run. An instrumented replay reproduces all
27 saved state tensors exactly and supplies the missing training curves.

On a newly acquired IES archive, the unchanged model correctly classified
27/27 agreed phase/background references at their centers and 26/27 after a
$+0.2$ s shift. Two prespecified primary cohorts together scored 44/44 at
center across 21 dates. Including three supplemental array-context references
produced 46/47 (97.87\%) at center and 45/47 (95.74\%) at $+0.2$ s. These are
conditional clip-classification results based on independently obtained,
model-blind AI visual references, not analyst-certified arrival labels or
continuous-stream detection accuracy. Limited sample size and unresolved
phase slots prevent a population-wide 95\% guarantee.

A C++17 implementation embeds the model constants, reads SAC or miniSEED and
returns three same-length probability vectors. All 141 reference-window
predictions match Python, with maximum absolute probability error
$4.17\times10^{-7}$. GitHub CI builds, tests and publishes Linux x86\_64/ARM64, macOS and
Windows executables under tag v1.0.0. A separate post-release sensitivity study completes 60 fits on date-grouped
training-pool folds, comparing two, three and four convolutional layers and
five QSIS epoch budgets. Additional depth shows no robust timing advantage;
longer fine-tuning gives modest validation gains while retaining a
fitting/validation gap. All figures in this report are PDF files.
\end{abstract}

\begin{table}[htbp]\centering\small
\caption{Principal findings and their distinct evaluation scopes.}
\begin{tabularx}{\linewidth}{Yrr}
\toprule Evaluation & Correct / scored & Accuracy \\\midrule
Prepared STEAD archive test, initial model & 16,426 / 17,449 & 94.14\% \\
Prepared QSIS test, initial fine-tuning & 176 / 248 & 70.97\% \\
Fresh IES confirmation, released model & 27 / 27 & 100.00\% \\
Combined primary confirmation, released model & 44 / 44 & 100.00\% \\
Primary plus supplemental references, center & 46 / 47 & 97.87\% \\
Primary plus supplemental references, $+0.2$ s & 45 / 47 & 95.74\% \\\bottomrule
\end{tabularx}
\end{table}
\clearpage
{\footnotesize\linespread{1.0}\selectfont\setlength{\parskip}{0pt}\tableofcontents}
\clearpage
{\small\setlength{\parskip}{0pt}\listoffigures}
\clearpage
\pagenumbering{arabic}

\section{Introduction and project scope}
\subsection{Motivation}
Dense, inexpensive three-component sensors can support the observation of
building response and local earthquake shaking. QSIS recordings are a useful
application domain because an indoor MEMS sensor sees not only regional
seismic motion but also structural amplification, changes in installation
orientation, quantization, transient interference and channel faults. A model
trained on a general seismic archive cannot be assumed to transfer unchanged
to these conditions. Prior work on phase detection and tinyML demonstrates
why compact learned waveform representations are relevant to this setting
\cite{10.1785/0120180080,10.1785/0220220322}.

The concrete objective here is three-class recognition of an already selected
short window. The project also provides a sliding-window interface for
continuous recordings. It does not establish event association, earthquake
location, warning lead time, structural damage assessment or an operational
false-alarm rate. Those applications require additional labels and downstream
validation. In particular, assigning a high P probability to a window does not
by itself establish the arrival time of its first P sample.

\subsection{Questions addressed in the restart}
The restart addressed four implementation concerns: model dimensions, use of
GPU training, lack of code organization and an incomplete C++ inference path.
The verified architecture has three convolutional and three fully connected
layers. The numbers 16, 32 and 64 are \emph{output channel counts}; every
convolution uses a $3\times1$ kernel. They should not be confused with kernel
sizes or temporal pooling dimensions. SpectrogramCNN is outside the retained
workflow. Training uses CUDA when available, otherwise Apple MPS, otherwise
CPU; the retained QSIS checkpoint was trained on MPS.

The active Python code is now organized into model, data, training, inference
and QSIS review modules. Historical scripts and artifacts remain available
under \code{legacy/} and \code{reports/legacy/}. The native implementation lives
under \code{cpp/}; the old Xcode prototype has been preserved separately at
\code{legacy/cpp/}. This report documents the retained model rather than treating
every exploratory experiment as a production change.

\section{Datasets, provenance and reference construction}
\subsection{Prepared archives and the initial baseline}
STEAD is a public collection of three-component earthquake and noise signals
\cite{8871127}. The local experiment uses prepared arrays derived from that
resource, not the complete global dataset. Shape checks, finite-value checks,
one-hot target validation, joint amplitude normalization and duplicate
screening precede fitting. Identical normalized training inputs are collapsed;
groups containing inconsistent labels are removed. The archived arrays lack
full source-event metadata, so the initial STEAD split is not an independently
verified event-disjoint split of the full upstream archive.

The local QSIS array contains adjacent P/S windows from each event recording
followed by noise windows. The baseline groups paired P/S windows by recording
before splitting. This prevents the same recording's P and S examples from
being assigned to separate partitions, but does not guarantee that different
stations recording one earthquake stay in the same partition. This baseline
is therefore distinct from the later event/date-separated confirmation work.

\begin{table}[htbp]\centering\small
\caption{Actual prepared-window counts in the saved baseline metrics, after
loading and filtering. These supersede unsourced counts in the old draft.}
\begin{tabularx}{\linewidth}{Yrrr}
\toprule Stage & Fitting & Validation & Test \\\midrule
Prepared STEAD, class-stratified development split & 58,740 & 14,686 & 17,449 \\
Prepared QSIS, recording-grouped split & 1,160 & 249 & 248 \\\bottomrule
\end{tabularx}
\end{table}

\subsection{Taiwan strong-motion adaptation}
The CWA Benchmark resource supplies TSMIP acceleration recordings at 200 Hz
and trace metadata \cite{cwaData}. The retained adaptation experiment uses the
2011 chunk and matching published sample metadata. It checks the component
order (ZNE), sampling rate, acceleration measurement, completeness flag,
single-event status and the metadata's manual-pick status, including the
published misspelling \code{maunal}. Accepted traces must be finite, have
nonconstant axes and leave sufficient context around the P and S references.

Exactly 2,748 traces survive. Each supplies P, S and pre-P background windows,
for 8,244 examples, balanced at 2,748 per class. Raw axes are reordered to ENZ.
The adaptation path detrends and tapers the complete trace, high-pass filters
at 0.5 Hz at its native rate and downsamples to 100 Hz before producing model
windows. This historical CWA preparation differs from the final QSIS
steady-start route; it is not silently described as identical preprocessing.

A seeded event split assigns 292 earthquake IDs to fitting and 73 to
validation. The resulting partitions contain 6,507 and 1,737 clips,
respectively. The chosen adaptation checkpoint is epoch 6, with validation
balanced accuracy 97.99\%. It is an initialization resource, and its validation
score is not a QSIS test result. The 2011 source period precedes the later QSIS
evaluation recordings.

\subsection{Final QSIS fitting pool}
The released checkpoint uses 226 ordered source windows: 33 P, 47 S and 146
noise windows. This pool combines the earlier 219-source recipe and seven
additional agreed CCUL training phases. It includes visually reviewed phases,
screened training proposals and background examples selected on training
dates. Proposals serve as weak training targets and are not promoted to
independent evaluation labels. Every source in the receipt records its paths,
center, target and group identity. The pool contains 63 distinct group
identifiers; event identifiers and noise-date identifiers have different
meanings and should not be interpreted as 63 independent earthquakes.

The predominance of noise windows motivates reciprocal-frequency class
weights. Per-record gravity alignment is usable for 104 fitting windows; the
remaining 122 retain the acquired axes. There is no station calibration table
in the released receipt. Training examples are not supplemented with any of
the fresh IES test clips. Figure~\ref{fig:composition} shows the class counts
and the exact composition of the confirmation references.
\fig[Fitting and confirmation class counts]{dataset_composition}{1}{Class composition of fitting and confirmation
sources. The two primary confirmation cohorts and three supplemental clips
are disjoint clip sets; they must not be counted a second time as separate
independent tests after aggregation.}{fig:composition}

\subsection{Prespecified primary confirmation}
The earlier primary confirmation cohort contains 17 scored references across
nine UTC dates: 3 P, 5 S and 9 pre-event backgrounds. Station/date enrollment
and model identities were fixed before target waveform review and scoring.
Nearby P-Alert recordings provide context, while the reference time is read
from the target QSIS waveform. Two independent model-blind AI reviewers
supply the waveform judgments. Of 36 enrolled phase slots across 18 dates,
eight yield agreed usable phase references. The unresolved slots are retained
in the coverage accounting rather than assigned convenient targets.

Additional array context later supports three more references on the reserved
dates: two P and one S. These are reported as a supplemental extension of the
17-clip set, giving 20 earlier references in total. Their review procedure is
still model-blind, but they are separated from the original primary enrollment
to preserve the distinction between prespecified and follow-up evidence.

\subsection{Fresh IES archive confirmation}
The IES Building Seismic Response Dataset was published on Zenodo in October
2025 and combines legacy SMA-1 and QSIS recordings \cite{iesData}. The local
archive inventory contains 3,069 QSIS SAC components, 112 event identifiers
and 64 event dates. The download is 225,569,391 bytes and matches the
publisher's MD5, \code{d0e9eef37f0844e349623174984ca8ea}. The original files
remain unchanged. The source supplies waveform/event metadata, but no usable
QSIS P/S analyst-arrival file was found in it.

Fifteen dates remain after excluding the original split, all candidate fitting
dates and all earlier public evaluation cohorts. Before visual inspection,
enrollment selects the largest listed earthquake per date, the two lowest
complete instrument codes as targets and up to two additional instruments as
context. Inputs must be aligned 100 Hz XYZ triplets with at least 30 seconds
before origin and 120 seconds afterward; nearby previous events are excluded
using a 120-second embargo. Catalog/SAC clock consistency is checked. Missing
station coordinates are supplied from the archive's documented IES location,
with that substitution recorded.

There are 29 enrolled target recordings, hence 58 possible P/S phase slots.
One target fails raw waveform QC, leaving 28 rendered targets. Two independent
AI readers examine the target waveforms in three bands, local array energy
and a wider overview without seeing classifier outputs. Agreement within
0.6 s yields nine P and six S references across six event dates. The other
41 rendered phase slots remain unresolved; the failed recording contributes
two additional QC-excluded slots. This is 15 agreed phase slots out of 58
enrolled, not an assumed 58/58 fully labelled benchmark.

Backgrounds are selected separately. A CWA catalog screen removes candidate
centers near regional catalog origins, and the maximum eligible 1 s / 5 s
STA/LTA peak is chosen in a prespecified pre-event interval. Raw dropouts and
inadequate filter/context support are excluded. Two readers accept 12 of 14
rendered background candidates, leaving a staircase-like quantized signal and
a disputed transient unresolved. The final fresh cohort is 9 P, 6 S and 12
noise clips on 12 dates. Figure~\ref{fig:coverage} makes the selection coverage
explicit.
\fig[Fresh IES reference coverage and exclusions]{reference_coverage}{.95}{Fresh reference coverage. The P and S rows each
contain 29 enrolled slots; one recording excludes one slot of each phase.
Accuracy is computed on the agreed references, not on unresolved slots.}{fig:coverage}
\FloatBarrier

\section{Waveform preprocessing and input representation}
\subsection{Synchronization and acquired-axis QC}
The verified Python route expects synchronized 100 Hz acceleration traces in
X, Y, Z order. All components need the same start, rate and sample count.
Missing or nonfinite data are rejected by the single-window Python interface.
The dense native interface instead splits finite segments and fills
unscorable output positions with NaN.

A center $c$ at 100 Hz uses samples $[c-100,c+100)$, giving exactly 200 raw
samples per component. A one-second repeated run in any \emph{acquired} axis
invalidates the window. This test precedes physical rotation: rotating a
constant or failed channel could otherwise mix motion from healthy axes into
it and conceal the fault. A zero or nonfinite normalization peak is also
invalid. Full two-second support is mandatory.

\subsection{Detrending and the steady-start high-pass filter}
Each complete finite recording is linearly detrended in double precision.
For a component $x_i$, the fitted line is $a+bi$, where $a,b$ minimize the
squared residual. A fourth-order Butterworth high-pass filter at 0.5 Hz is
applied as two second-order sections. Instead of tapering the beginning of a
recording, the released QSIS path uses steady initial conditions obtained
from \code{sosfilt_zi} multiplied by the first detrended sample. This avoids a
taper amplitude ramp while preserving the selected training/inference route.
The filter is causal, but whole-record detrending uses later samples. The
current file-based pipeline is therefore \emph{not} a strictly online causal
picker; streaming deployment would require a separately validated detrending
and state-handling design.

\subsection{Gravity-derived rotation}
Stable gravitational offsets can identify the vertical axis of a stationary
acceleration sensor. The algorithm computes one-second component medians over
seconds 1--7, then a componentwise median across those six blocks. Let that
baseline be $\mathbf{g}$ and its normalized direction be
$\mathbf{u}=\mathbf{g}/\|\mathbf{g}\|$. Alignment requires a magnitude between
900 and 1,100 gal and a maximum block-vector drift of at most 8 gal, where
$1\;\mathrm{gal}=1\;\mathrm{cm\,s^{-2}}$. The shortest proper rotation maps
$\mathbf{u}$ to the positive vertical axis. If the baseline is unusable, the
identity transform is retained.

For $\mathbf{v}=\mathbf{u}\times\mathbf{e}_z$ and
$d=\mathbf{u}\cdot\mathbf{e}_z$, the nondegenerate rotation can be written
\begin{equation}
 R=I+[\mathbf{v}]_\times+
 \frac{[\mathbf{v}]_\times^2}{1+d},
\end{equation}
with explicit parallel/antiparallel handling. Alignment is applied only from
center sample 800 onward, so the six-second baseline precedes the complete
classification window. It determines vertical orientation, not geographic
azimuth; it cannot establish north/east directions from gravity alone. There
is no station-name fallback in the deployed model.

\subsection{Fourier resampling and joint normalization}
After filtering and, when applicable, rotation, a 200-sample cut is
Fourier-resampled to 80 samples per axis. At 40 Hz this remains a two-second
window. All components share one peak scale:
\begin{equation}
 \widetilde{x}_{t,a}=
 \frac{x_{t,a}}{\max_{t',a'}|x_{t',a'}|},
 \qquad t=0,\ldots,79,\quad a\in\{X,Y,Z\}.
\end{equation}
The tensor supplied to PyTorch has shape $[B,1,80,3]$ in float32. Joint
normalization preserves relative component amplitudes, although it removes
absolute amplitude. Figure~\ref{fig:examples} shows both acquired deviations
and the actual normalized model inputs for representative held-out clips.

No time-frequency image is supplied to this classifier.
Figure~\ref{fig:spectral} is an explanatory diagnostic calculated from the
same reference waveforms; it does not change the inference representation.
The class spectra describe this small selected sample and do not establish a
universal spectral boundary between phases or between buildings.
\fig[Raw and normalized held-out waveform examples]{waveform_examples}{1}{Representative fresh held-out P, S and noise
windows. Upper panels subtract the per-axis median only for display; lower
panels show the released pipeline's normalized waveform input. X/Y/Z are
colored consistently. The captioned softmax values come from the frozen
checkpoint.}{fig:examples}
\fig[Class spectra and diagnostic time-frequency view]{spectral_diagnostics}{1}{Waveform diagnostics, not model inputs. Left:
summed three-axis Welch power from the 80-sample normalized inputs, with
classwise median and interquartile range across the 47 centered references.
Right: an STFT of filtered vector magnitude around an example P reference;
its color is log amplitude and its time scale differs from the two-second
classification cut.}{fig:spectral}
\FloatBarrier

\section{CNN architecture}
Figure~\ref{fig:architecture} summarizes the retained processing chain.
Convolutions have unit stride, no padding and $3\times1$ kernels. Temporal
max pooling has kernel and stride $3\times1$. Thus early convolutions share
temporal filters across the three physical components without convolving
across the component dimension. The flattened fully connected stage combines
those component-specific feature responses.
\fig[Preprocessing and retained CNN architecture]{pipeline_architecture}{1}{Released preprocessing and CNN. Output
channels are 16/32/64; temporal pooling reduces 80 model samples to two time
positions. The final feature volume is $64\times2\times3=384$.}{fig:architecture}

\begin{table}[htbp]\centering\small
\caption{Layer dimensions and trainable parameters; batch dimension omitted.
BN parameters include learned scale/bias, not running-statistic buffers.}
\begin{tabularx}{\linewidth}{Yrr}
\toprule Operation & Output shape & Parameters \\\midrule
Input waveform & $1\times80\times3$ & 0 \\
Conv1, $1\to16$, kernel $3\times1$ & $16\times78\times3$ & 64 \\
BN1, GELU, max pool, dropout 0.1 & $16\times26\times3$ & 32 \\
Conv2, $16\to32$, kernel $3\times1$ & $32\times24\times3$ & 1,568 \\
BN2, GELU, max pool, dropout 0.1 & $32\times8\times3$ & 64 \\
Conv3, $32\to64$, kernel $3\times1$ & $64\times6\times3$ & 6,208 \\
BN3, GELU, max pool, dropout 0.1 & $64\times2\times3$ & 128 \\
Flatten & $384$ & 0 \\
FC1, GELU & $64$ & 24,640 \\
FC2, GELU & $16$ & 1,040 \\
FC3, logits & $3$ & 51 \\\midrule
Total & & \textbf{33,795} \\\bottomrule
\end{tabularx}
\end{table}

The exact GELU is $x\Phi(x)$, with $\Phi$ the standard normal cumulative
function. Softmax maps logits $z_k$ to
$p_k=\exp(z_k-m)/\sum_j\exp(z_j-m)$, $m=\max_jz_j$. Class order is always
P, S, Noise, corresponding to indices 0, 1, 2. Dropout is disabled and BN
uses stored running statistics at evaluation. A longer-input research path
contains adaptive pooling, but the selected checkpoint and native port use
exactly 80 input samples and need no adaptive-pooling operation.
\FloatBarrier

\section{Training recipe and recorded learning curves}
\subsection{Initial prepared-archive training}
The saved initial STEAD run uses MPS, PyTorch 2.11.0, seed 42, batch size 256,
50 epochs and Adam with learning rate $10^{-3}$. StepLR halves the learning
rate every seven epochs. Validation selects epoch 25, rather than assuming
the final epoch is best. The held-out prepared archive then scores
16,426/17,449. Figure~\ref{fig:stead} reproduces the original epoch log.

The initial QSIS archive fine-tuning starts from that STEAD checkpoint,
uses AdamW, batch size 32, 100 epochs, initial learning rate $10^{-3}$ and
the same seven-epoch schedule. Validation selects epoch 10. The test score is
176/248, with substantial P/S confusion. Figure~\ref{fig:qsis-old} shows that
training accuracy keeps improving while validation largely plateaus. This
is evidence of a generalization gap, not a reason to report the training
accuracy as achieved QSIS accuracy. Its test partition is left separate from
the later reviewed references and their differently defined labels.
\fig[Prepared STEAD training and validation curves]{training_stead}{1}{Original 50-epoch prepared STEAD training log. The
vertical line marks the validation-selected checkpoint at epoch 25. Test
metrics are not evaluated at every epoch or drawn as a selection curve.}{fig:stead}
\fig[Initial QSIS training and validation curves]{training_qsis_baseline}{1}{Original 100-epoch QSIS archive fine-tuning.
The selected checkpoint is epoch 10. The fitting/validation gap and persistent
phase confusion motivate the subsequent review and adaptation work.}{fig:qsis-old}

\subsection{CWA adaptation}
The CWA stage initializes from the STEAD checkpoint, updates all CNN weights
using AdamW at $10^{-4}$, batch size 128, six epochs and seed 42. It applies
time-shift and horizontal-rotation augmentation to training batches. The
class balance is exact, so ordinary cross-entropy is used. Validation is
held out by earthquake identifier. All six recorded epochs are shown in
Figure~\ref{fig:cwa}; the retained epoch-6 checkpoint initializes the final
QSIS fit. Its checksum is recorded in Appendix~\ref{app:provenance}.
\fig[CWA adaptation training and validation curves]{training_cwa}{.97}{Six recorded CWA adaptation epochs. Validation is
held out by CWA event ID and measures the adaptation domain. It is not the
independent QSIS confirmation set.}{fig:cwa}

\subsection{Released QSIS fine-tuning}
The final fit updates all layers for a fixed 20 epochs with seed 42 on MPS.
Batch size is 16; AdamW uses learning rate $3\times10^{-4}$ with its default
weight decay 0.01, betas $(0.9,0.999)$ and epsilon $10^{-8}$. There is no
learning-rate scheduler in this final fit. The last epoch is retained; the
fresh test is not used to select an epoch.

Training augmentation reflects time indices at the clip boundaries rather
than wrapping them. Each example receives an integer shift from $-8$ to $+8$
at 40 Hz (up to $\pm0.2$ s), and a horizontal rotation uniformly drawn from
$\pm30^\circ$. The vertical axis remains unchanged, and joint peak
normalization is reapplied. These transformed inputs are training views of
the same 226 source windows; they do not increase the number of independent
source recordings.

Reciprocal-frequency class weights $w_k=1/n_k$ enter weighted cross-entropy:
\begin{equation}
 L_{\mathcal{B}}=
 -\frac{\sum_{i\in\mathcal{B}}w_{y_i}\log p_{i,y_i}}
         {\sum_{i\in\mathcal{B}}w_{y_i}}.
\end{equation}
The existing epoch logger weights these batch means by batch size before
averaging. Because class weights and batch compositions vary, that logged
training objective is not mathematically identical to one weighted reduction
over the entire epoch. The clean fitting-set diagnostic uses ordinary
unweighted cross-entropy. Figure~\ref{fig:final-fit} labels the two loss
reductions explicitly; their numerical gap should not be interpreted solely
as overfitting.

\subsection{Recovery and verification of the missing final curves}
The original final fitting function saved weights and source receipts but no
per-epoch history. To recover truthful curves, the reporting utility reloads
the 226 source paths/centers in their recorded order and reruns the existing
fitting function with the same initialization, seed, device, augmentation and
optimizer. A wrapper records each existing epoch result and evaluates only
the clean \emph{fitting} inputs in evaluation mode. CPU and MPS random states
are preserved around the diagnostic. No validation or test reference is read
by this replay.

After fitting, every state tensor is compared with the released checkpoint
using exact tensor equality. All 27 tensors match. The curves are therefore
an instrumented exact replay of the released weights, not fabricated or
smoothed estimates of an unavailable training log. The clean fitting score
at epoch 20 is 224/226 (99.12\%); the online augmented training score is
222/226 (98.23\%). These remain training scores. A high fitting score alone
cannot demonstrate freedom from overfitting.
\fig[Exact replay of the final 20-epoch fit]{training_final}{1}{Exact 20-epoch replay of the released final fit.
Blue curves are online augmented training metrics with dropout active;
orange curves evaluate the unaugmented fitting inputs after each epoch with
dropout disabled and stored BN statistics. There is no test/validation curve
for this final refit. All 27 saved state tensors reproduce exactly.}{fig:final-fit}

\subsection{What constrains overfitting}
The project uses transfer initialization, a small model, dropout, weight
decay, class weighting, training augmentation, event/date exclusions and
separate confirmation cohorts. Exploratory changes are evaluated on
development data. Larger networks are not adopted merely because their
training loss improves. Nonetheless, many development comparisons have been
made, and archived evaluation references were subsequently reused. They are
reported descriptively, rather than retrospectively relabelled as untouched
holdouts. The fresh cohort is the clearest additional confirmation because
both model identities and its enrollment protocol were fixed before its
waveform review.
\FloatBarrier

\section{Sensitivity to model depth and training duration}\label{sec:sensitivity}
\subsection{Study design and separation from confirmation}
A post-release sensitivity study varies convolutional depth and the number
of QSIS fitting epochs. Its data are restricted to the 226 fitting sources
already recorded in the released receipt. None of the 47 independent
confirmation references is used to select an architecture, an epoch or an
initialization. This is development cross-validation on the existing fitting
pool, whose targets include reviewed and weak training labels; its scores
are not new independent confirmation accuracy.

To avoid recording/date leakage, source-event dates, SAC start dates and
center-sample dates are joined into connected UTC-date groups. This also
keeps an earthquake crossing midnight from being split through its start and
origin dates. The 226 sources form 58 such groups. An exact audit also finds
226 unique normalized centered inputs and no identical input shared across
validation folds. Five folds are frozen
before fitting. Their validation class counts (P/S/Noise) are $(7,10,28)$,
$(6,10,28)$, $(7,9,30)$, $(7,9,28)$ and $(6,9,32)$. Fold assignment is balanced
using target counts only, without inspecting predictions. Each validation
source is held out once per seed. A group's sources cannot appear on both
sides of its fold.

Three fixed fitting seeds, 42, 43 and 44, are paired across all configurations.
Every trajectory runs for 80 epochs, saving evaluations at epochs 5, 10, 20,
40 and 80. All configurations use the same source cuts, five folds, batch
size 16, AdamW learning rate $3\times10^{-4}$, weight decay 0.01, reciprocal
\emph{fitting-fold} class weights and original shift/rotation augmentation.
Raw evaluation windows are recut at $-0.2$, 0 and $+0.2$ s. Validation uses
evaluation mode and unweighted cross-entropy. No early stopping or adaptive
rerun is triggered by these validation measurements.

The study completes $4\times3\times5=60$ fits and 4,800 aggregate QSIS
optimization epochs, or 57,600 optimizer steps. This is the experiment-wide
budget, not the age of a single model. Predictions from all five validation
folds are pooled over 226 sources separately for each seed. Tables and plots
then report the mean and sample standard deviation across the three seeds.
The five folds and repeated timing views are not counted as independent
replicates to narrow an error bar. Seed SD is not a confidence interval.

\subsection{Depth variants and a fair initialization comparison}
The depth comparison keeps the fully connected classifier fixed at
$384\rightarrow64\rightarrow16\rightarrow3$ and compares two, three and four
convolutional layers, hence five, six and seven learned Conv/FC layers in
that counting convention. BN affine operations are not counted as extra
Conv/FC layers. All three depth variants start from random weights: the
pretrained three-layer checkpoint cannot initialize a different-width
architecture without a transfer confound.

\begin{table}[htbp]\centering\small
\caption{Depth variants. Pooling differences are explicit; this is an
architecture sensitivity study rather than a claim to change only an
abstract scalar depth while holding every feature operation identical.}
\begin{tabularx}{\linewidth}{lrrY}
\toprule Conv / FC layers & Channels & Parameters & Temporal reduction \\\midrule
2 / 3 & 16/64 & 29,091 & Two original $3\times1$ max pools leave eight time positions; adaptive average pooling reduces them to two before flattening. \\
3 / 3 & 16/32/64 & 33,795 & Exact original architecture, with three $3\times1$ max pools. \\
4 / 3 & 16/32/64/64 & 46,275 & A padded $3\times1$ 64-channel block is inserted before the final max pool; all three temporal downsamplings are retained. \\\bottomrule
\end{tabularx}
\end{table}

A fourth configuration uses the exact three-convolution model with the
original CWA initialization. It tests epoch duration under the released
transfer recipe without implying that random and pretrained starts are the
same condition. No alternative model is promoted into the native binary.
\fig[Depth sensitivity across grouped folds]{sensitivity_depth}{1}{Depth
sensitivity from random initialization. Each cell is a mean over three seeds
of pooled five-fold out-of-fold predictions. The right panel takes the
minimum balanced accuracy across the three fixed raw-window offsets within
each seed before averaging. The physical source count remains 226.}{fig:sensitivity-depth}

\subsection{Depth results}
At 20 epochs, centered balanced accuracy is $91.63\pm1.69$\% for two
convolutions, $93.29\pm1.19$\% for three and $92.86\pm0.78$\% for four.
At 80 epochs, the corresponding results are $93.65\pm0.67$\%,
$93.75\pm1.58$\% and $94.01\pm0.96$\%. The four-layer mean is only 0.26
percentage points above the three-layer mean at epoch 80, while using
36.93\% more parameters. That difference is small relative to the observed
seed variation and does not establish a dependable benefit of greater depth.

The two-layer variant has 13.92\% fewer parameters than the original and
similar centered balanced accuracy at epoch 80. However, the models' timing
robustness differs: their mean worst-offset balanced accuracies at epoch 80
are 91.86\%, 91.54\% and 90.33\%, respectively. The deeper variant's slight
centered gain therefore does not persist under the timing perturbation.
Figure~\ref{fig:sensitivity-capacity} compares all trajectories and capacity.
\fig[Model capacity and validation performance]{sensitivity_capacity}{1}{All
four development trajectories and their parameter counts. Each marker in
the capacity panel is one prespecified epoch checkpoint; several markers
share one parameter count. This is not an independent test-set ranking.}{fig:sensitivity-capacity}

\subsection{Epoch sensitivity of the released architecture}
For the CWA-initialized original CNN, the mean centered accuracy rises from
93.07\% at five epochs to 95.72\% at 20 and 96.90\% at 80. Centered balanced
accuracy is $90.03\pm0.92$\%, $92.90\pm0.47$\% and $93.83\pm1.01$\%,
respectively. Thus the larger 20-to-80-epoch budget improves mean centered
accuracy by 1.18 percentage points but balanced accuracy by only 0.93 points.
The fitting-fold accuracy reaches 100\% at epoch 80, while out-of-fold
accuracy remains 96.90\%; fitting accuracy alone would conceal this gap.
Mean centered P recall is 87.88\% at both 20 and 80 epochs; S recall rises
from 92.20\% to 93.62\% and noise recall from 98.63\% to 100\%. The overall
accuracy gain therefore does not imply that the harder P examples improve.
The class-recall panel preserves this distinction.

Worst-offset balanced accuracy is 91.22\% at 20 epochs, 91.14\% at 40 and
91.84\% at 80. The effect is not a simple monotonic improvement at every
checkpoint. The original random-start CNN's balanced accuracy also drops
from 93.29\% at epoch 20 to 92.41\% at 40 before recovering to 93.75\% at 80.
These results argue for measuring grouped validation and timing robustness
rather than assuming either more epochs or more layers is automatically
better. The full grid, including all intermediate epochs, is in
Appendix~\ref{app:sensitivity-grid}.
\fig[Epoch sensitivity and the generalization gap]{sensitivity_epochs}{1}{Epoch
sensitivity of the exact original architecture under random and CWA starts.
Upper left: mean centered out-of-fold balanced accuracy with seed-SD bars.
Upper right: clean fitting-fold and pooled out-of-fold accuracy. Lower left:
mean worst-offset balanced accuracy. Lower right: mean centered class recall
for the CWA start. The independent confirmation set is not evaluated along
these trajectories.}{fig:sensitivity-epochs}

\subsection{Meaning of total epochs and the retained model}
The varied epoch count is the final QSIS fitting duration, with preceding
initialization held fixed to isolate that factor. The CWA start descends from
STEAD's saved epoch 25 and six CWA adaptation epochs. Its saved-weight
ancestry therefore contains $25+6+E$ stage epochs: 36, 41, 51, 71 and 111
for the five values of $E$. The original STEAD job actually executed 50
epochs while selecting epoch 25; including that discarded work gives
$50+6+E$ executed pipeline epochs. Neither sum is a common unit of compute,
because the stages have different dataset sizes and batch sizes. The earlier
100-epoch prepared-QSIS baseline is a separate experiment, not an ancestor
of this CWA-initialized checkpoint.

The retained full-data model remains the original 20-epoch checkpoint, and
its SHA-256 is checked unchanged after this study. The cross-validation
improvement at 80 epochs is a useful development observation, not permission
to re-label the existing 47 clips as a new untouched test for an optimized
model. Any replacement should first be frozen and then confirmed on new
expert-labelled events and background periods. The 58-group training-pool
study does not remove its weak-label or single-network limitations.
\FloatBarrier

\section{Evaluation protocol and statistics}
\subsection{Units, offsets and metrics}
The evaluation unit is one two-second reference clip. Its target is the
phase/noise category of a selected window, not a dense sequence of analyst
labels at every input sample. Each clip is re-cut from the raw recording at
center offsets $-0.2$, 0 and $+0.2$ s. Those three views are strongly related;
47 clips produce 141 evaluations, not 141 independent test examples.

Confusion matrices use reference classes as rows and predicted classes as
columns, in the order P, S, Noise. For matrix $C$, class recall is
$C_{kk}/\sum_j C_{kj}$ and precision is $C_{kk}/\sum_i C_{ik}$. Accuracy is
$\operatorname{tr}(C)/\sum_{ij} C_{ij}$; balanced accuracy averages class recalls,
and macro F1 averages the three class F1 scores. These definitions make the
minority P class visible even when noise dominates the overall score.

Report generation reopens each of the 47 references, applies the released
preprocessing and recomputes all 141 softmax vectors using the frozen model.
Every argmax prediction matches the saved comparison. The regenerated
matrices are also checked against the immutable evaluation receipt, so the
figures and tables use one consistent prediction source.

\subsection{Reference independence and coverage}
Two independently obtained model-blind AI reviews are separate waveform
judgments, but they are not seismologist-certified manual arrival files.
Agreement does not eliminate correlated visual interpretation errors. No
expert QSIS arrival catalog was supplied. References from unclear targets
are excluded from accuracy and counted explicitly, as in
Figure~\ref{fig:coverage}. The 0.6 s reviewer-agreement tolerance should not
be read as achieved arrival-time error by the classifier.

The combined primary sample has 44 clips across 21 dates; the three
supplemental references reuse those reserved dates and bring the total to 47.
Many phase clips share an event, building or recording. The reference counts
must therefore be accompanied by date/site coverage, not treated as a large
collection of independent random draws from all future QSIS observations.

\subsection{Uncertainty and the 95\% objective}
For transparency, Figure~\ref{fig:timing} includes ordinary 95\% Wilson
intervals treating clips as independent Bernoulli observations. For 46/47,
that illustrative interval is 88.89--99.62\%; for 44/44 it is approximately
91.97--100\%. Independence is an approximation because clips cluster by
UTC date and site. These intervals already show why the observed 97.87\%
point estimate is not a lower-bound guarantee above 95\%.

A second descriptive analysis resamples all clips jointly by UTC date,
using 10,000 seeded replicates over the 21 dates. It gives percentile
intervals 93.02--100\% at center and 89.19--100\% at $+0.2$ s for the
47-clip set. Resampling the observed dates cannot introduce unseen buildings,
unresolved reference types or unobserved failures. All-correct subsets can
also give degenerate bootstrap intervals. Neither calculation removes the
selection and label-coverage limitations.

Even under an ideal independent-binomial interpretation, 59 consecutive
correct trials would be required for a one-sided exact 95\% lower bound to
exceed 95\%, since $0.05^{1/59}>0.95$. A realistic QSIS study also needs
independent events, sites and noise conditions. This numerical illustration
is a study-design guide, not a claim that adding a few selected clips would
resolve the current evidence limitations.
\FloatBarrier

\section{Results}
\subsection{Centered confusion matrices}
Figure~\ref{fig:confusion} presents all three centered confirmation scopes.
The fresh IES matrix is diagonal with counts $(9,6,12)$. The combined primary
matrix is diagonal with counts $(12,11,21)$. In the 47-clip set, one of the
14 reference P clips is classified as S, giving
\begin{equation}
 C_{47,0}=\begin{pmatrix}13&1&0\\0&12&0\\0&0&21\end{pmatrix}.
\end{equation}
The single centered error remains in every reported aggregate. None of the
supplemental references is moved back into fitting after inspection.
\fig[Centered confirmation confusion matrices]{confusion_matrix_test}{1}{Centered confirmation confusion matrices.
Cells show counts and row-normalized percentages. Left: fresh IES only.
Middle: the two original primary cohorts. Right: primary plus three
supplemental array-context references. The panels overlap by construction
and must not be summed.}{fig:confusion}

\begin{table}[htbp]\centering\small
\caption{Centered per-class performance on the 47-clip extended confirmation
set. All precision/recall values are conditional on scored references.}
\begin{tabular}{lrrrr}
\toprule Reference class & Support & Precision & Recall & F1 \\\midrule
P & 14 & 100.00\% & 92.86\% & 96.30\% \\
S & 12 & 92.31\% & 100.00\% & 96.00\% \\
Noise & 21 & 100.00\% & 100.00\% & 100.00\% \\\midrule
Macro average & 47 & 97.44\% & 97.62\% & 97.43\% \\\bottomrule
\end{tabular}
\end{table}

\subsection{Timing sensitivity}
At $-0.2$ s the extended matrix is the same as at center. At $+0.2$ s,
one additional P reference becomes S:
\begin{equation}
 C_{47,+0.2}=\begin{pmatrix}12&2&0\\0&12&0\\0&0&21\end{pmatrix}.
\end{equation}
The resulting 45/47 (95.74\%) still exceeds the numerical 95\% point-estimate
target on the scored references. P recall at this offset falls to 12/14
(85.71\%), so aggregate accuracy should not conceal a weaker class-specific
result. The fresh cohort alone has 26/27, with 8/9 P correct.
\fig[Extended confirmation at fixed timing offsets]{confusion_matrix_offsets}{1}{Extended confirmation matrices at the
three fixed offsets. The two errors at $+0.2$ s are both P-to-S confusions;
all scored S and noise windows remain correct.}{fig:offsets}
\fig[Timing sensitivity and illustrative uncertainty]{timing_and_uncertainty}{1}{Timing sensitivity and sample-size context.
The left panel shows fixed-offset accuracy for overlapping confirmation
scopes. The right panel shows illustrative clip-iid Wilson intervals at
center; dependence and selected-reference coverage prevent their use as a
population-wide performance guarantee.}{fig:timing}

\subsection{Baseline test results are retained}
The prepared STEAD and early QSIS test matrices are shown in
Figure~\ref{fig:old-matrices}. Their denominators and label construction
differ from the later confirmation cohorts. The early QSIS matrix has
P recall 41/72, S recall 38/72 and noise recall 97/104; its balanced accuracy
is 67.66\%. Its 70.97\% accuracy is a concrete domain-transfer warning.
The revised 97.87\% result should not be described as a controlled
26.90-percentage-point improvement on an identical test set. Both model
training and reference definition changed, and the datasets are different.
\fig[Historical prepared-archive test matrices]{confusion_matrix_historical}{.98}{Saved initial prepared-archive test
results. Rows are reference classes and cells contain count plus row
percentage. These are historical baselines, not the final model's
independent confirmation denominators.}{fig:old-matrices}

\subsection{Failure analysis and probability diagnostics}
The centered miss is the supplemental CCUL P reference GPH-015 on
2025-03-25. The retained gravity model predicts S at all three offsets.
At IES reference IHP-017 on 2024-04-17, the centered P decision changes to S
at $+0.2$ s. Figure~\ref{fig:errors} shows their local filtered waveforms and
all three class probabilities. These examples are preserved as errors;
they are not corrected by changing labels or by tuning a post-hoc threshold
on the confirmation cohort.

The waveform/probability views help identify overlap and timing sensitivity,
but they do not establish a unique physical cause for the errors. A dense
S probability near an already known P reference could reflect later motion
inside the same two-second cut, sensor/site transfer or a model limitation.
Analyst review and independent follow-up data are needed to separate these
possibilities. Probability magnitudes are shown in
Figure~\ref{fig:confidence}, without claiming calibration from only 47
selected clips.
\fig[Two retained P-to-S failures]{error_cases}{1}{The two distinct extended-confirmation P-to-S errors.
Upper panels show eight seconds of filtered XYZ context and shade the
centered two-second cut. Lower panels use freshly recomputed frozen-model
softmax probabilities at the three fixed offsets.}{fig:errors}
\fig[Reference probabilities and native numerical parity]{confidence_and_native_parity}{1}{Left: probability assigned to each
centered reference class, with the centered error ringed. This is not a
reliability/calibration estimate. Right: native-versus-Python maximum
absolute probability differences on 141 related window evaluations; all
argmax decisions agree.}{fig:confidence}

\subsection{Direct waveform-to-inference comparison}\label{sec:continuous-display}
To make the model behavior visible along a recording, the native executable
processes the complete held-out IHP-003 XYZ triplet from IES station 5AE21.
The recording contains 30,001 samples at 100 Hz. Dense CPU inference uses
stride one, produces 29,802 finite probability triples and leaves 199
unsupported edge centers as NaN. It is run on the entire recording before
selecting a display interval, so the plot does not accidentally change
whole-record detrending or filter initialization through a shorter recut.
At both reviewed reference centers, the native probabilities agree with the
Python values to within $2\times10^{-6}$.

Figure~\ref{fig:continuous-event} aligns the three acquired axes and class
probabilities on one time coordinate. A fixed componentwise baseline is
removed only from the displayed raw waveform; the inference path still uses
the released preprocessing. Dashed lines mark the model-blind P and S
reference times, 22.85 and 38.25 s after catalog origin, and shaded bands mark
their two-second reference windows. Figure~\ref{fig:continuous-full} retains
the complete approximately 300-second view, including the pre-event interval
and later coda.

The plotted probabilities are center-window scores, not per-sample onset
labels. In particular, later shaking alternates between high S and high
noise scores, and short P-score excursions also occur away from the reviewed
P reference. The rest of the recording has no dense expert target sequence,
so these plots neither establish continuous detection accuracy nor measure
false alarms per hour. They provide the direct waveform/inference comparison
missing from the earlier clip-only figures.
\fig[Event waveform and dense class probabilities]{waveform_inference_event}{1}{Waveform
and dense inference on the same shared time axis, with an 82-second display
excerpt chosen from the full-record run. Top three panels: acquired XYZ
acceleration minus a fixed display baseline. Bottom: unsmoothed P/S/Noise
probabilities at every supported 100 Hz center. Dashed lines and shaded
windows mark the two reviewed reference centers. The displayed sample
probabilities are not automatically interpreted as arrival picks.}{fig:continuous-event}
\fig[Full-record waveform and inference comparison]{waveform_inference_full_record}{1}{The
same dense inference over the entire 30,001-sample held-out recording.
The 199 unsupported edge centers are NaN, without artificial Noise labels.
Only the two reviewed phase windows have reference categories; the complete
curve is a qualitative diagnostic rather than a continuously labelled test.}{fig:continuous-full}

\subsection{Archived reuse and background stress evidence}
A separate 215-clip archived reference set scores 210/215 at center for the
current model, with 211/215 at $-0.2$ s and 209/215 at $+0.2$ s. At center,
its matrix is $((47,3,0),(2,65,0),(0,0,98))$. These references had already
been used in development comparisons and are descriptive reuse, not a new
untouched test. They are not added to the 47-clip confirmation denominator.

The retained model also correctly classifies 35/35 archived trigger-noise
clips at center, with 34/35 at the negative offset. On 40 reviewed interior
noise references, date-held-out fitting gives 40/40 across all three offsets
at seeds 42 and 43. These are useful targeted background checks with their
own protocols. They do not establish a false-alarm count per hour on an
uninterrupted noise stream, and their overlapping histories are not pooled
into a single headline percentage.
\FloatBarrier

\section{Model comparisons and 2026 research context}
\subsection{Local comparisons and retained choices}
The development work considered capacity, context and training-set changes,
not just additional epochs. Table~\ref{tab:experiments} summarizes selected
comparisons already recorded in the benchmark notes. Different rows have
different matched subsets; their percentages must not be ranked as if they
came from one common test set. They also precede the fresh IES confirmation.

\begin{table}[htbp]\centering\small
\caption{Selected exploratory comparisons. These are development or archived
reuse measurements and are not additional independent final tests.}
\label{tab:experiments}
\begin{tabularx}{\linewidth}{>{\raggedright\arraybackslash}p{.27\linewidth}>{\raggedright\arraybackslash}p{.22\linewidth}Y}
\toprule Experiment & Matched result & Implication \\\midrule
Augmented two-second CNN, seeds 42--44 & 42, 42, 41 / 46 & Small seed variation remains; timing shifts matter. \\
Polarization residual head & 41/46 vs CNN 42/46 & No centered benefit in the matched trial. \\
Four-second context CNN & 40/46 vs 42/46 & Longer context alone did not improve development recognition. \\
Real-noise mixing, 8--20 dB & 42, 40, 42 / 46 vs 42, 42, 41 & No consistent gain across the matched seeds. \\
Frozen U-Trans encoder + regularized head & 35/41 vs CNN 36/41 & Five clips lack 60 s context; this limited adaptation does not favor replacement. \\
234-source expanded CNN & Same fresh 27/27; 26/27 at $+0.2$ s & Earlier S/noise regressions favor retaining the 226-source model. \\
Equal probability ensemble & No centered gain over expanded component on 52 phases & More inference cost without consistent benefit. \\\bottomrule
\end{tabularx}
\end{table}

A smaller 5,955-parameter TCN, smoothed temporal pooling, multiscale residual
features and optional history-based refinement remain research alternatives
in the repository. The released native program implements the selected
single CNN and its recorded preprocessing; it does not silently add a
probability ensemble or the optional history head. That choice keeps model
identity and native parity unambiguous.

\subsection{Recent primary literature}
PhaseNet established a dense three-component P/S arrival-probability approach
\cite{phasenet}. Its task differs from assigning one label to a selected
window. This distinction also applies to newer networks and prevents direct
comparison of a picking F1 score with QSIS clip accuracy.

DTPP (2026) combines depthwise separable temporal processing and dilations
for efficient phase picking \cite{dtpp}. Its relevance here is a compact
longer-context design principle. The local TCN experiment is inspired by
that general direction; it is not a replication of DTPP and the retained
CNN is not renamed as that architecture.

U-Trans (March 2026) pretrains a U-Net/convolutional-transformer waveform
representation by reconstructing corrupted inputs \cite{utrans}. It is a
reasonable candidate when transferable context representations are needed.
Our local comparison freezes its backbone and fits a regularized head on
small QSIS folds. The published framework also studies joint downstream
fine-tuning, so the local 35/41 comparison is evidence about this restricted
configuration, not a general refutation of foundation models.

PhaseNeXt was published on 4 November 2026 and studies a neural picker trained
on twenty years of JMA records \cite{phasenext}. Its ConvNeXt/multiscale context
and use of carefully checked small-event examples are relevant to future
QSIS timing work. A local timing pilot used the authors' pinned weights:
27 of 46 previously reviewed public arrivals passed the raw/band-pass
agreement filter, and 26 of those were within 0.6 s of the visual reference.
That is a proposal-quality diagnostic with incomplete retention, not 26/27
end-to-end QSIS phase classification accuracy. The earlier RED-PAN search
matched 36 references and the union matched 40; differing search windows
and preprocessing prevent an architecture-only comparison.

These sources support continued investigation of transferable context and
label quality. They do not establish that replacing a 33,795-parameter CNN
with a transformer will improve this small sensor-specific dataset. A useful
next comparison would freeze a new event/site test protocol, hold out expert
arrivals, match context availability and measure class recall, timing error,
continuous false alarms and CPU cost together.
\FloatBarrier

\section{Standalone C++ inference and reproducible release}
\subsection{Public API and input decoding}
The C++17 project exposes standard-library types through
\code{cpp/include/seismictorch/inference.hh}. Three binary SAC v6/v7 files
are accepted in X/Y/Z order, including either byte order and the SAC v7
double footer. Component rates, UTC starts and sample counts must agree.
A whole miniSEED 2 or 3 file is decoded by a statically linked vendored
libmseed 3.6.0. Compressed Steim and supported numeric encodings are handled
without requiring a separately installed decoder at runtime.

miniSEED components are grouped by network, station, location and channel
family; XYZ, ENZ and 12Z suffixes map to the three axes. An ambiguous
multi-station file requires an explicit stream selector. Samples are aligned
to a common UTC grid. Gaps become NaN, while conflicting overlapping samples
are rejected. Files are limited to 2 GB and the aligned grid to 50 million
samples. Inputs need calibrated acceleration; the gravity check expects gal.
Instrument response removal and automatic unit inference are not performed.

\begin{lstlisting}[language=C++]
#include <seismictorch/inference.hh>

auto wave = seismictorch::read_sac({"X.sac", "Y.sac", "Z.sac"});
// Or: auto wave = seismictorch::read_mseed("recording.mseed");
auto scores = seismictorch::infer(wave);
std::vector<float> P = std::move(scores.p);
std::vector<float> S = std::move(scores.s);
std::vector<float> Noise = std::move(scores.noise);
\end{lstlisting}

Every vector has exactly the aligned input length. The default stride is one
sample, so every supported sample center is evaluated. The first/last
unsupported centers, nonfinite gaps, raw dropout windows and zero signals
remain NaN in all three vectors. They are not fabricated noise predictions.
Stride values above one explicitly request an approximation with interpolation
between valid anchors; they are not identical to dense inference.

\subsection{Raspberry Pi 5 / 64-bit ARM Linux}
The CI matrix now includes \code{linux-arm64} on the native GitHub
\code{ubuntu-24.04-arm} runner. Its Alpine/musl container builds with baseline
\code{-march=armv8-a} C/C++ flags, static standard runtimes and the same
embedded model. CI executes all four CTest checks natively on ARM64, verifies
the ELF AArch64 machine type and absence of a dynamic interpreter, then runs
an actual 1,400-sample miniSEED inference outside the build container.
This is execution verification, not only cross-compilation.

The target is Raspberry Pi 5 with 64-bit Raspberry Pi OS or another 64-bit
Linux distribution. No matching glibc installation is required. A 32-bit OS
cannot execute this package. Hosted ARM64 testing does not provide a
measurement of physical Pi 5 throughput, power use or thermal behavior.
The ARM64 archive is added to the existing v1.0.0 release together with an
updated four-package checksum file; the original three archives and tag are
retained. Future tagged CI releases automatically require all four builds.
The native ARM64 evidence is
\href{https://github.com/Gauss963/SeismicTorch/actions/runs/44905643313}{run 44905643313}.

\subsection{Embedded constants and numerical equivalence}
The Python exporter reads the verified checkpoint and training receipt,
checks the expected CNN/preprocessing and writes hexadecimal floating-point
constants to \code{cpp/generated/model_weights.hh}. Batch normalization is
folded into each convolution:
\begin{equation}
 W'_o=\frac{\gamma_o}{\sqrt{\sigma_o^2+\epsilon}}W_o,\qquad
 b'_o=\frac{\gamma_o(b_o-\mu_o)}{\sqrt{\sigma_o^2+\epsilon}}+\beta_o.
\end{equation}
The generated constants are committed, so normal CMake compilation requires
neither Python nor a model file. Convolutions, exact GELU, pooling, fully
connected layers and stable softmax execute directly on CPU. The native
program does not require PyTorch, LibTorch, BLAS or MLX at runtime.

The original verified preprocessing path is 100 Hz. Native rates 40--1000 Hz
are accepted by constructing the high-pass filter at the acquired rate and
resampling each two-second window to 80 model samples. Golden preprocessing
fixtures check 40, 50, 100, 200 and 1000 Hz; this is numerical implementation
verification, not an accuracy benchmark for every sensor/rate combination.

Full local parity uses 27 actual SAC recordings containing the 47 references,
recomputed at all three offsets. The native path includes SAC decoding,
detrending, high-pass initialization, gravity alignment, Fourier resampling
and the embedded CNN. All 141 decisions match Python; the maximum absolute
probability difference is $4.1723164\times10^{-7}$. Smaller CI fixtures carry
the same normalized model windows and independent preprocessing golden cases,
without transferring the full local waveform archives to every build runner.

\subsection{Cross-platform CMake and GitHub release}
Owned C++ sources consistently use \code{.cc} and \code{.hh}. Unmodified
third-party C retains its upstream \code{.c}/\code{.h} extensions and license.
CMake builds the static library target \code{SeismicTorch::seismictorch} and
the executable \code{seismictorch-infer}. External CMake consumers can use
\code{add_subdirectory} and link that alias. A separate consumer build was
verified to obtain three vectors of length 1,400 from a miniSEED fixture.

\begin{table}[htbp]\centering\small
\caption{Published v1.0.0 native packages and actual dependency checks.}
\begin{tabularx}{\linewidth}{lY}
\toprule Platform & Build and dependency policy \\\midrule
Ubuntu x86\_64 & Fully static musl executable; no dynamic interpreter, and execution checked outside the build container. \\
macOS universal & arm64 and x86\_64 slices, deployment target 11.0; only system libSystem/libc++ dependencies. \\
Linux ARM64 / Pi 5 & Fully static ARMv8-A musl executable; four tests pass on a native ARM64 Linux runner. Requires a 64-bit OS. \\
Windows x86\_64 & Static Microsoft C/C++ runtime; dependency inspection reports only KERNEL32.dll. \\\bottomrule
\end{tabularx}
\end{table}

Four CTest checks pass on each platform: native model/reader/preprocessing
checks, command help, embedded model/version identity and a full CLI
roundtrip. The CLI test requires exactly 1,400 data rows and identical CSV
output for equivalent SAC and miniSEED inputs. Reader tests also cover
endianness, SAC v7, miniSEED 2/3, compressed samples, CRC rejection, Unicode
paths, multiple stream families, gaps, overlaps, raw dropout and vector
lengths. The release job waits for all four builds, packages license notices,
generates SHA-256 checksums and publishes tagged releases. The original
v1.0.0 assets passed the three-platform pipeline; a later native ARM64 run
adds the fourth architecture without changing the compiled inference sources
or embedded model. Downloaded archives were
checksum-verified, and the downloaded macOS executable passed universal-slice,
signature and inference-output checks.

\begin{lstlisting}
cmake -S cpp -B build/native -DCMAKE_BUILD_TYPE=Release
cmake --build build/native --config Release --parallel
ctest --test-dir build/native -C Release --output-on-failure

seismictorch-infer --sac X.sac Y.sac Z.sac --output scores.csv
seismictorch-infer --mseed recording.mseed --output scores.csv
\end{lstlisting}

The CSV columns are \code{time_seconds,P,S,Noise}. Time is relative to the
aligned recording start. Finite entries sum to one; unsupported entries are
written as \code{nan}. Although the output is dense, its probabilities still
represent the sliding two-second classifier. Whole-record detrending,
centered windows and gravity initialization prevent interpreting the current
program as a zero-latency real-time alert system.
\FloatBarrier

\section{Limitations and next validation steps}
The retained model meets the numerical 95\% objective on the scored,
independently reviewed confirmation references at all three fixed offsets.
That statement has a defined denominator and should remain attached to it.
The following limits materially affect interpretation:
\begin{enumerate}
\item \textbf{Reference quality.} No analyst-certified QSIS arrivals were
available. AI visual agreement is useful audit evidence but cannot replace
external seismological annotation for a deployment claim.
\item \textbf{Partial phase coverage.} Only 15 of 58 newly enrolled phase
slots are agreed/scored. Excluding unclear examples makes the task
conditional on readability; unresolved examples cannot be assumed correct.
\item \textbf{Sample size and dependence.} Forty-seven clips on 21 dates
remain small and correlated. The illustrative uncertainty bounds extend below
95\%; all-correct subsets provide little information about rare failures.
\item \textbf{Domain coverage.} The new archive emphasizes magnitude-5-or-
larger earthquakes at one building. Other buildings, weaker arrivals,
instrument changes and long quiet periods need independent coverage.
\item \textbf{Historical reuse.} The 215 archived references and many
exploratory comparisons have informed development. Their scores are not
fresh final-test evidence.
\item \textbf{Task mismatch.} Clip accuracy does not measure pick timing,
event-level detection, missed-event rates or false alarms per hour. A
continuous stream contains many overlapping windows and many more background
opportunities than the confirmation clips.
\item \textbf{Temporal availability.} Centered windows require later samples,
gravity estimation needs a stable preceding baseline, and whole-record
linear detrending is offline. A causal runtime would need its own training
and validation protocol.
\item \textbf{Probability interpretation.} Softmax values have not been
independently calibrated. A high maximum value is not a guaranteed
probability of correctness under a new sensor distribution.
\end{enumerate}

The next study should obtain expert references, freeze event/site-disjoint
training/validation/test groups, document all unscorable intervals and
annotate continuous background hours. Thresholds and debounce logic should
be selected on validation only. The final test should report per-class
recall/F1, event-level sensitivity, matched-arrival timing error, false alarms
per hour, CPU throughput and unavailable-output coverage. A transformer or
multiscale replacement should be compared on exactly those same frozen
inputs and resource constraints. Additional training should be motivated by
those failures rather than tuning against the present confirmation cases.

\section{Conclusion}
The restart turns the original prototype into a reproducible waveform
classification and native inference project. The retained architecture is
verified as a 16/32/64-channel CNN with a 384/64/16/3 classifier, trained on
Apple MPS and reproduced exactly after code reorganization. The report
preserves the weak early QSIS baseline and documents how stronger
arrival-reference review, regional adaptation and preprocessing discipline
led to a final 46/47 centered extended-confirmation score. The 45/47 positive-
offset result and two P-to-S failures remain visible.

The standalone v1.0.0 release embeds the selected weights and provides tested
SAC/miniSEED inference on Linux, macOS and Windows. Numerical agreement with
Python is established on all 141 fixed-offset evaluations. The date-grouped sensitivity study does not support increasing depth merely
for a small centered gain. An 80-epoch QSIS fit is a development candidate
for further confirmation, while the released 20-epoch weights remain fixed.
The next research
priority is a larger expert-labelled, event/site-separated and continuous
QSIS benchmark, so the observed clip-level result can be tested against the
broader reliability requirements of real deployment.

\clearpage
\appendix
\section{Reproduction and artifact guide}\label{app:reproduce}
\subsection{Report build}
Run the figure generator in the project's Python environment. All source
snapshots needed to redraw the report are stored under
\code{reports/legacy/final-report/report-data/}. The ordinary figure build
therefore does not require downloading the full local experimental archives.
All newly generated figures are vector PDF; raster PNGs used for layout
inspection are temporary and are not manuscript figures.
\begin{lstlisting}
.venv/qsis-review/bin/python \
  reports/legacy/final-report/tools/make_figures.py
.venv/qsis-review/bin/python \
  reports/legacy/final-report/tools/make_extended_figures.py
cd reports/legacy/final-report
tectonic --keep-logs Project_Report.tex
\end{lstlisting}
\code{tectonic} runs the bibliography and required cross-reference passes.
A conventional LaTeX/BibTeX installation can also build the same source.
The output is \code{Project_Report.pdf}; the original title/logo, author and
advisor information are preserved. Scientific text, measurements and figures
have been rewritten for the current implementation.

\subsection{Refresh measurements from the local waveform evidence}
The collection script reopens the original local SAC files, records their
hashes, recomputes all class probabilities, checks saved predictions/matrices
and snapshots the required histories. The fitting replay is separate and
requires the original MPS environment and source waveforms. Neither script
replaces the released model.
\begin{lstlisting}
.venv/qsis-review/bin/python \
  reports/legacy/final-report/tools/collect_evidence.py
.venv/qsis-review/bin/python \
  reports/legacy/final-report/tools/reproduce_training.py
\end{lstlisting}

The complete 226-source training recipe is documented in
\code{docs/qsis-benchmark.md}; it uses the frozen split, reviewed training
arrivals, the training-only proposal/trigger pools, CWA initialization,
20 epochs, seed 42, augmentation, steady-start filtering and gravity alignment.
The reproduction script uses the selected receipt to recover those exact
ordered windows and verify the final tensors, avoiding an undocumented new
fit on a different data subset.

\subsection{Snapshot contents}
\begin{longtable}{p{.35\linewidth}p{.59\linewidth}}
\toprule Artifact in report-data & Purpose \\\midrule\endhead
\code{evidence.json} & Saved baseline metrics, epoch histories, final receipts, cohort comparisons, C++ parity and release evidence. \\
\code{test-predictions.json} & Forty-seven raw-path references, waveform hashes, all three probability vectors, checked matrices, class metrics and date-bootstrap settings. \\
\code{sensitivity.json} & Frozen date-grouped protocol, 60 completed fits, all fold/seed predictions, full 20-row depth/epoch grid and initialization scope. \\
\code{continuous-event.npz} & Full 30,001-sample acquired waveform and native dense probabilities used in the aligned display. \\
\code{arm64-ci.json} & Successful four-architecture CI job results; ARM package identity and checksum are in \code{arm64-package.json}. \\
\code{final-training-replay.json} & Twenty recovered epochs, initialization/input/checkpoint hashes, device/version and exact state-tensor equality result. \\
\code{waveform-panels.npz} & Frozen raw deviations/context inputs and normalized cuts used to redraw waveform/spectral figures. \\
\code{provenance.json} & Original source paths and SHA-256 checksums, class/orientation counts and matrix conventions. \\
\bottomrule
\end{longtable}

\clearpage
\section{Complete sensitivity grid}\label{app:sensitivity-grid}
All values below are percentages. Accuracy, balanced accuracy (BA), and
worst-offset BA are pooled out of fold within each seed, then averaged over
seeds 42--44. The spread after BA is the sample SD over those three seeds.
These 20 rows are the complete prespecified grid rather than selected
favorable checkpoints. Epochs count final QSIS fitting updates. Every model
keeps three fully connected layers.
\begingroup\small
\begin{longtable}{rrrrrrr}
\caption{Five-fold date-grouped development sensitivity on 226 fitting sources.}\\
\toprule Conv & Init & Epoch & Parameters & Accuracy & BA $\pm$ SD & Worst BA \\\midrule\endfirsthead
\toprule Conv & Init & Epoch & Parameters & Accuracy & BA $\pm$ SD & Worst BA \\\midrule\endhead
\bottomrule\endfoot
2 & Random & 5 & 29,091 & 91.30 & 86.03 $\pm$ 0.68 & 83.52 \\
2 & Random & 10 & 29,091 & 92.63 & 88.54 $\pm$ 1.02 & 85.73 \\
2 & Random & 20 & 29,091 & 94.69 & 91.63 $\pm$ 1.69 & 90.07 \\
2 & Random & 40 & 29,091 & 94.99 & 92.88 $\pm$ 0.60 & 91.28 \\
2 & Random & 80 & 29,091 & 96.17 & 93.65 $\pm$ 0.67 & 91.86 \\
3 & Random & 5 & 33,795 & 92.18 & 86.02 $\pm$ 3.56 & 83.68 \\
3 & Random & 10 & 33,795 & 92.63 & 88.80 $\pm$ 1.47 & 86.68 \\
3 & Random & 20 & 33,795 & 95.58 & 93.29 $\pm$ 1.19 & 91.31 \\
3 & Random & 40 & 33,795 & 94.69 & 92.41 $\pm$ 1.28 & 90.60 \\
3 & Random & 80 & 33,795 & 96.17 & 93.75 $\pm$ 1.58 & 91.54 \\
4 & Random & 5 & 46,275 & 88.94 & 86.27 $\pm$ 0.76 & 81.85 \\
4 & Random & 10 & 46,275 & 91.30 & 89.36 $\pm$ 1.26 & 86.40 \\
4 & Random & 20 & 46,275 & 94.54 & 92.86 $\pm$ 0.78 & 90.24 \\
4 & Random & 40 & 46,275 & 95.28 & 92.93 $\pm$ 0.25 & 90.21 \\
4 & Random & 80 & 46,275 & 95.43 & 94.01 $\pm$ 0.96 & 90.33 \\
3 & CWA & 5 & 33,795 & 93.07 & 90.03 $\pm$ 0.92 & 86.85 \\
3 & CWA & 10 & 33,795 & 95.43 & 93.07 $\pm$ 0.90 & 90.75 \\
3 & CWA & 20 & 33,795 & 95.72 & 92.90 $\pm$ 0.47 & 91.22 \\
3 & CWA & 40 & 33,795 & 96.46 & 93.24 $\pm$ 0.83 & 91.14 \\
3 & CWA & 80 & 33,795 & 96.90 & 93.83 $\pm$ 1.01 & 91.84 \\

\end{longtable}
\endgroup

\section{Model identity and source integrity}\label{app:provenance}
The released model is \code{models/qsis/current/model.pth}, with SHA-256:
\begin{quote}\small\nolinkurl{634eea951ff2f8e4a42f4bbf12174aaa0a8aa6699d15726d4cba9e8cf90fe424}\end{quote}
Its immutable fitting receipt has SHA-256:
\begin{quote}\small\nolinkurl{0e9449932f21ff2248b4a183b303fe845132cbcc2d41e65ad856e9b3f9a5d283}\end{quote}
The CWA adaptation initializer has SHA-256:
\begin{quote}\small\nolinkurl{6911d989b0052d06de5aa9afad47f450b92e3729f22e6551e9631379d95625ea}\end{quote}
The STEAD pretraining checkpoint has SHA-256:
\begin{quote}\small\nolinkurl{ecb6502114f64d8730ced85de3f411bc74062e314616604d1b10299a900ccb34}\end{quote}

The model's original fitting receipt retains the status string
\code{development_candidate_unvalidated}, because it was written before later
confirmation. The subsequent evaluation receipts, rather than a retroactive
edit to that training artifact, record the new reviewed evidence. The source
snapshot/hashes and exact migration replay are recorded in
\code{docs/reorganization-verification.json}. Its historical C++-prototype
scope describes the earlier reorganization stage; the later native v1.0.0
implementation is verified separately.

The native release points to source commit
\code{6f65971f677abc887ad37192799ccc646733c032}. The complete successful tagged
CI run is \href{https://github.com/Gauss963/SeismicTorch/actions/runs/44903391622}{GitHub Actions run 44903391622}.
The published archives and their checksums are at
\url{https://github.com/Gauss963/SeismicTorch/releases/tag/v1.0.0}.

\section{Every scored confirmation reference}
Table~\ref{tab:allrefs} lists each distinct reference once. P and S from one
recording share an identifier but are separate labelled clips. The first
20 rows are the earlier primary/supplemental cohort; the following 27 are
fresh IES references. Rows marked by predictions differing from the reference
are retained in the confusion matrices. Exact raw paths, sample centers,
source categories and three-way probabilities are in
\code{test-predictions.json}.
\begingroup\small
\begin{longtable}{lllllll}
\caption{All 47 confirmation clips. Prediction columns are fixed center offsets
in seconds.}\label{tab:allrefs}\\
\toprule ID & Reference & UTC date & Site & $-0.2$ & $0$ & $+0.2$ \\\midrule\endfirsthead
\toprule ID & Reference & UTC date & Site & $-0.2$ & $0$ & $+0.2$ \\\midrule\endhead
\midrule\multicolumn{7}{r}{Continued on next page}\\\endfoot
\bottomrule\endlastfoot
GPH-001 & P & 2023-10-22 & RCEC & P & P & P \\
GPH-001 & S & 2023-10-22 & RCEC & S & S & S \\
GPH-015 & S & 2025-03-25 & CCUL & S & S & S \\
GPH-016 & P & 2025-08-24 & HW02 & P & P & P \\
GPH-016 & S & 2025-08-24 & HW02 & S & S & S \\
GPH-017 & P & 2025-08-27 & HW02 & P & P & P \\
GPH-017 & S & 2025-08-27 & HW02 & S & S & S \\
GPH-018 & S & 2025-12-23 & CHY01 & S & S & S \\
GPB-001 & Noise & 2023-10-22 & RCEC & Noise & Noise & Noise \\
GPB-011 & Noise & 2024-05-09 & CHY01 & Noise & Noise & Noise \\
GPB-012 & Noise & 2024-08-31 & Fy0 & Noise & Noise & Noise \\
GPB-013 & Noise & 2024-10-24 & RCEC & Noise & Noise & Noise \\
GPB-014 & Noise & 2024-12-11 & Fy0 & Noise & Noise & Noise \\
GPB-015 & Noise & 2025-03-25 & CCUL & Noise & Noise & Noise \\
GPB-016 & Noise & 2025-08-24 & HW02 & Noise & Noise & Noise \\
GPB-017 & Noise & 2025-08-27 & HW02 & Noise & Noise & Noise \\
GPB-018 & Noise & 2025-12-23 & CHY01 & Noise & Noise & Noise \\
GPH-012 & P & 2024-08-31 & Fy0 & P & P & P \\
GPH-013 & S & 2024-10-24 & RCEC & S & S & S \\
GPH-015 & P & 2025-03-25 & CCUL & S & S & S \\
IHP-003 & P & 2022-10-10 & IES & P & P & P \\
IHP-003 & S & 2022-10-10 & IES & S & S & S \\
IHP-004 & P & 2022-10-10 & IES & P & P & P \\
IHP-004 & S & 2022-10-10 & IES & S & S & S \\
IHP-005 & S & 2022-10-28 & IES & S & S & S \\
IHP-006 & S & 2022-10-28 & IES & S & S & S \\
IHP-007 & P & 2022-11-01 & IES & P & P & P \\
IHP-008 & P & 2022-11-01 & IES & P & P & P \\
IHP-009 & P & 2022-11-20 & IES & P & P & P \\
IHP-009 & S & 2022-11-20 & IES & S & S & S \\
IHP-010 & P & 2022-11-20 & IES & P & P & P \\
IHP-010 & S & 2022-11-20 & IES & S & S & S \\
IHP-012 & P & 2023-10-14 & IES & P & P & P \\
IHP-013 & P & 2023-10-14 & IES & P & P & P \\
IHP-017 & P & 2024-04-17 & IES & P & P & S \\
IBG-001 & Noise & 2021-12-30 & IES & Noise & Noise & Noise \\
IBG-003 & Noise & 2022-10-10 & IES & Noise & Noise & Noise \\
IBG-005 & Noise & 2022-10-28 & IES & Noise & Noise & Noise \\
IBG-007 & Noise & 2022-11-01 & IES & Noise & Noise & Noise \\
IBG-009 & Noise & 2022-11-20 & IES & Noise & Noise & Noise \\
IBG-011 & Noise & 2022-12-15 & IES & Noise & Noise & Noise \\
IBG-012 & Noise & 2023-10-14 & IES & Noise & Noise & Noise \\
IBG-016 & Noise & 2024-04-17 & IES & Noise & Noise & Noise \\
IBG-018 & Noise & 2024-04-23 & IES & Noise & Noise & Noise \\
IBG-022 & Noise & 2024-06-10 & IES & Noise & Noise & Noise \\
IBG-024 & Noise & 2024-06-29 & IES & Noise & Noise & Noise \\
IBG-028 & Noise & 2024-11-30 & IES & Noise & Noise & Noise \\

\end{longtable}
\endgroup

\FloatBarrier
\clearpage
\bibliographystyle{unsrt}
\bibliography{bibliography}
\end{document}
